% =====================================================================
% Chapter 4: 3-point local and logarithm terms, Eqs. (I.54)-(I.67)
% Tasks 9 and 10 of the re-derivation of Paper I.
% =====================================================================
\providecommand{\chIVloc}{\mathrm{loc}}
\providecommand{\chIVlog}{\mathrm{log}}
\providecommand{\chIVok}{\ensuremath{\checkmark}}
\providecommand{\chIVno}{\ensuremath{\times}}

\chapter{3-point local and logarithm terms: Eqs. (I.54)--(I.67)}
\label{ch:ch04}

This chapter covers two consecutive tasks of the re-derivation of \PI.
Paper~I splits the 3-point contribution to the scalar self-mass into three parts:
\emph{local} contributions, which come from the delta-function term of $\partial_0\partial'_0 i\Delta_A$;
\emph{logarithm} contributions, which come from the de Sitter breaking factor $k\ln(aa')$ of the outer propagator $i\Delta_A$;
and \emph{normal} contributions, which contain everything else.
Task~9 (Section~\ref{sec:ch04:task9}) re-derives the local part, Eqs.~\eqI{54}--\eqI{59}.
Task~10 (Section~\ref{sec:ch04:task10}) re-derives the logarithm part, Eqs.~\eqI{60}--\eqI{67}.
The normal contributions, \eqI{68} and later, are treated in later chapters.

\begin{notebox}[Notation and conventions used in this chapter]
\begin{itemize}[leftmargin=1.2em,itemsep=1pt]
\item $k\equiv\dfrac{H^{D-2}}{(4\pi)^{D/2}}\dfrac{\Gamma(D-1)}{\Gamma(\frac D2)}$ (Paper~I, Eq.~\eqI{13}); as $D\to4$, $k\to\dfrac{H^2}{8\pi^2}$.
\item $\pihat\equiv-\pi\cot\big(\tfrac{\pi D}{2}\big)=-\dfrac{2}{D-4}+\order{D-4}$.
\item $p\equiv D-2$; $\;l\equiv\ln a(\eta)$, $\;l'\equiv\ln a(\eta')$, $\;L\equiv\ln(aa')=l+l'$.
Since $a=-1/(H\eta)$, we have $\partial_0a=Ha^2$ and therefore $\partial_0 l=Ha$.
\item Labels: \tP\ given in the paper, \tD\ derived here, \tA\ assumption or convention, \tC\ check by code.
\item \textbf{Sign convention.} Every expression \emph{derived} in this chapter is in the \emph{definition sign}.
This is the sign fixed by Paper~I's own definition \eqI{168}, $\partial_\mu(\sqrt{-g}\,g^{\mu\nu}\partial_\nu\Phi)-\int M^2\Phi=0$ with $-iM^2=\delta^2 i\Gamma$.
It is the sign obtained from the derived \eqI{47}--\eqI{49} of Task~7.
Every self-mass contribution \emph{printed} in Paper~I is in the \emph{paper's convention}, which is $-1$ times the definition sign.
(Identities that do not involve the self-mass, such as \eqI{55}--\eqI{56} and \eqI{60}--\eqI{62}, are convention-independent.)
\item \textbf{Input principle.} We use only results \emph{derived} in earlier tasks:
\eqI{47}--\eqI{49} (Task~7), the local term of \eqI{52} (Task~8), the coincidence limits (Task~4),
and the distributional lemma and the full-$D$ form of \eqI{44} (Task~5).
\end{itemize}
\end{notebox}

% =====================================================================
\section{Task 9: the local 3-point contributions, Eqs. (I.54)--(I.59)}
\label{sec:ch04:task9}
% =====================================================================

Task~9 derives the local 3-point contributions with the corrected sign and normalizations and compares them with the paper (printed page 124005-6).

\begin{discbox}[Summary of findings (Task 9)]
\begin{itemize}[leftmargin=1.2em,itemsep=2pt]
\item The derived \eqI{54}, \eqI{57} and \eqI{58} are exactly $-1$ times the printed ones.
On each line, the paper's right-hand side is a correct evaluation of its own left-hand side.
The left-hand sides, however, carry the signs of the \emph{printed} \eqI{47}--\eqI{49}.
So the paper carried the sign error of the step \eqI{46}$\to$\eqI{47}--\eqI{49} into the rest of the calculation.
This answers the question left open in Task~7: there it was undecided whether the later steps of the paper used the printed signs of \eqI{47}--\eqI{49} or the signs implied by \eqI{46}.
The first possibility holds.
\item \textbf{Arithmetic error in \eqI{59}.} Adding the paper's own \eqI{54}, \eqI{57} and \eqI{58} gives the finite coefficient $-\tfrac52$ for $a^{D-2}\nabla^2$.
The paper prints $-\tfrac72$.
The difference has the form of the counterterm \eqI{145} and can be absorbed with $\Delta\alpha_3=-\tfrac{1}{8\pi^2}$.
\item \textbf{Typo in the text.} The paper refers to ``the 4th term of (47)''.
As the left-hand side of \eqI{54} shows, the term meant is the 2nd term of \eqI{47}.
\item \textbf{Important correction (of Task~6).} In the paper, the 4-point and 3-point contributions are consistent with each other: both have the opposite sign relative to the definitions \eqI{168} and \eqI{131}.
Hence the paper's bare $M^2$ is $-1$ times the one defined by \eqI{168}.
\end{itemize}
\end{discbox}

% ---------------------------------------------------------------------
\subsection{Goal, inputs and scope}
% ---------------------------------------------------------------------

\paragraph{Target equations \tP.}
Paper~I (p.~124005-6, left column) prints the following.
It introduces \eqI{54} as the local contribution from ``the 4th term of (47)'' without the logarithm:
\begin{align*}
&-\kappa^2\Big(\frac{D-1}{D-3}\Big)\partial_0\partial'_0\Big[(aa')^{D-2}A(y)\,\frac{i}{a^{D-2}}\,\delta^D(x-x')\Big]\\
&\qquad=\frac{i\kappa^2H^{D-2}}{(4\pi)^{D/2}}\,\frac{\Gamma(D)}{(D-3)\Gamma(\frac D2)}\Big(-\pi\cot\big(\tfrac D2\pi\big)\Big)
\times\big\{-a^D\square\delta^D(x-x')+a^{D-2}\nabla^2\delta^D(x-x')\big\},
\tag{I.54}\\[4pt]
&-\partial_0\big(a^{D-2}\partial'_0\delta^D(x-x')\big)=-\partial^\mu\big(a^{D-2}\partial_\mu\delta^D(x-x')\big)+a^{D-2}\nabla^2\delta^D(x-x')
\tag{I.55}\\
&\phantom{-\partial_0\big(a^{D-2}\partial'_0\delta^D(x-x')\big)}\equiv-a^D\square\delta^D(x-x')+a^{D-2}\nabla^2\delta^D(x-x'),
\tag{I.56}\\[4pt]
&-\kappa^2\nabla\!\cdot\!\nabla'\Big[(aa')^{D-2}B(y)\times\frac{i}{a^{D-2}}\,\delta^D(x-x')\Big]
=-\frac{i\kappa^2H^2}{16\pi^2}\times a^2\nabla^2\delta^4(x-x')+\order{D-4},
\tag{I.57}\\[4pt]
&2\kappa^2\Big(\frac{D-2}{D-3}\Big)\partial_0\partial'_0\Big[(aa')^{D-2}C(y)\times\frac{i}{a^{D-2}}\,\delta^D(x-x')\Big]\\
&\qquad=\frac{i\kappa^2H^2}{4\pi^2}\big\{a^4\square\delta^4(x-x')-a^2\nabla^2\delta^4(x-x')\big\}+\order{D-4},
\tag{I.58}\\[4pt]
&-iM^2_{\rm 3pt,\,loc}(x;x')=\frac{i\kappa^2H^{D-2}}{(4\pi)^{D/2}}\frac{\Gamma(D-1)}{\Gamma(\frac D2)}
\times\Big\{\Big[\Big(\frac{D-1}{D-3}\Big)\pi\cot\big(\tfrac\pi2D\big)+2\Big]a^D\square\notag\\
&\qquad\qquad+\Big[-\Big(\frac{D-1}{D-3}\Big)\pi\cot\big(\tfrac\pi2D\big)-\frac72\Big]a^{D-2}\nabla^2+\order{D-4}\Big\}\delta^D(x-x').
\tag{I.59}
\end{align*}
In \eqI{54} the factor $-\pi\cot(\frac D2\pi)$ is printed after the fraction without parentheses; it multiplies the fraction.
None of these printed expressions was used as an input. They enter only the comparison in Section~\ref{sec:ch04:t9comp}.

\paragraph{Inputs.}
\begin{itemize}[leftmargin=1.2em,itemsep=3pt]
\item[(G1)] The derived \eqI{47}--\eqI{49} \tD\ (Task~7).
Only three of the nine terms have the inner derivative pair $\partial_0\partial'_0$, and only these give a local contribution (Task~8):
\begin{equation}
\begin{aligned}
\Tab_{AA}&=+\kappa^2\,\frac{D-1}{D-3}\,\partial_0\partial'_0\Big[(aa')^{D-2}\,i\Delta_A\,\partial_0\partial'_0\,i\Delta_A\Big],\\
\Tab_{AB}&=+\kappa^2\,\nabla\!\cdot\!\nabla'\Big[(aa')^{D-2}\,i\Delta_B\,\partial_0\partial'_0\,i\Delta_A\Big],\\
\Tab_{AC}&=-2\kappa^2\,\frac{D-2}{D-3}\,\partial_0\partial'_0\Big[(aa')^{D-2}\,i\Delta_C\,\partial_0\partial'_0\,i\Delta_A\Big].
\end{aligned}
\label{eq:ch04:G1t9}
\end{equation}
In each term the outer derivatives (in front of the square bracket) act on everything inside the bracket.
The inner derivatives act only on the second propagator, $i\Delta_A$.
In the paper's \eqI{47}--\eqI{49} these three terms appear with the opposite signs.
\item[(G2)] The local part of \eqI{52} \tD\ (Task~8):
\begin{equation}
\partial_0\partial'_0\,i\Delta_A(x;x')\supset\frac{i}{a^{D-2}}\,\delta^D(x-x').
\label{eq:ch04:G2t9}
\end{equation}
\item[(G3)] Coincidence limits \tD\ (Task~4):
\begin{equation}
A(x;x)=k\pihat,\qquad b\equiv i\Delta_B(x;x)=-\frac{k}{D-2},\qquad c\equiv i\Delta_C(x;x)=\frac{k}{(D-2)(D-3)}.
\label{eq:ch04:G3t9}
\end{equation}
\item[(G4)] Separation convention \tP.
For $I=A$, the paper uses only $A(y)$ in the local contribution.
The piece $k\ln(aa')\times\delta^D$ goes to the logarithm contributions \eqI{60}--\eqI{67}.
In Task~8 we had put this piece on the local side instead.
Both separations are complete; for the comparison we adopt the paper's separation.
\end{itemize}

\paragraph{Scope.}
The logarithm and normal contributions are not treated in this task.
$D$ is kept general.
Coefficients multiplying a pole are kept at exact $D$; $D=4$ is set only in finite parts that are not multiplied by a pole.

% ---------------------------------------------------------------------
\subsection{Rules \texorpdfstring{\tD}{}}
% ---------------------------------------------------------------------

\paragraph{(R1) Products with the delta function.}
For a smooth function $g$, the identity $g(x,x')\,\delta^D(x-x')=g(x,x)\,\delta^D(x-x')$ holds as a distribution.
Derivatives are taken \emph{after} this identity has been applied.
The equality $A(y)\,\delta^D(x-x')=k\pihat\,\delta^D(x-x')$ is an assumption of dimensional regularization \tA.
For $\mathrm{Re}\,D<2$, $A$ is continuous at coincidence and all powers of $y$ vanish there.
The result is then carried to general $D$ by analytic continuation.
With \eqref{eq:ch04:G2t9} this gives, for $I=A$ (without its logarithm), $B$ and $C$:
\begin{equation}
(aa')^{D-2}\,i\Delta_I(x;x')\,\frac{i}{a^{D-2}}\,\delta^D(x-x')=i\,\Lambda_I\,a^{D-2}\,\delta^D(x-x'),
\qquad \Lambda_A=k\pihat,\quad \Lambda_B=b,\quad \Lambda_C=c .
\label{eq:ch04:R1}
\end{equation}
Here $(aa')^{D-2}/a^{D-2}=a'^{D-2}$ becomes $a^{D-2}$ on the support of the delta function.

\paragraph{(R2) Outer derivatives.}
We act with the kernel on a test function $g$.
Since $a^{D-2}=a(\eta)^{D-2}$ depends only on $x$, the primed derivative acts only on the delta function:
\begin{equation}
\int d^Dx'\;\partial_\mu\partial'_\nu\big[a^{D-2}\delta^D(x-x')\big]\,g(x')=\partial_\mu\big[a^{D-2}\,(-\partial_\nu g)\big].
\end{equation}
By the definition \eqI{42}, $a^D\square=\partial^\mu(a^{D-2}\partial_\mu)$. With $\eta^{00}=-1$ this reads
$a^D\square f=-\partial_0(a^{D-2}\partial_0f)+a^{D-2}\nabla^2f$. Therefore
\begin{equation}
\partial_0\partial'_0\big[a^{D-2}\delta^D\big]=-\partial_0\big(a^{D-2}\partial_0\delta^D\big)=\big(a^D\square-a^{D-2}\nabla^2\big)\delta^D,
\qquad
\nabla\!\cdot\!\nabla'\big[a^{D-2}\delta^D\big]=-a^{D-2}\nabla^2\delta^D .
\label{eq:ch04:R2}
\end{equation}
The paper's \eqI{55}--\eqI{56} is the first equality of \eqref{eq:ch04:R2} multiplied by $-1$:
$-\partial_0(a^{D-2}\partial'_0\delta^D)=-a^D\square\delta^D+a^{D-2}\nabla^2\delta^D$.
This is correct~\chIVok.

% ---------------------------------------------------------------------
\subsection{The three local terms \texorpdfstring{\tD}{}}
% ---------------------------------------------------------------------

\paragraph{Eq.~\eqI{54}: $\Tab_{AA}$.}
We insert \eqref{eq:ch04:R1} with $\Lambda_A=k\pihat$ and then \eqref{eq:ch04:R2}.
Using $(D-1)\Gamma(D-1)=\Gamma(D)$:
\begin{equation*}
\begin{aligned}
\Tab_{AA}^{\chIVloc}&=\kappa^2\,\frac{D-1}{D-3}\cdot i\,k\pihat\,\big(a^D\square-a^{D-2}\nabla^2\big)\delta^D\\
&=\frac{i\kappa^2H^{D-2}}{(4\pi)^{D/2}}\,\frac{\Gamma(D)}{(D-3)\Gamma(\frac D2)}\Big(-\pi\cot\big(\tfrac{\pi D}{2}\big)\Big)
\big\{+a^D\square-a^{D-2}\nabla^2\big\}\delta^D(x-x').
\end{aligned}
\tag{I.54$'$}\label{eq:ch04:54}
\end{equation*}

\paragraph{Eq.~\eqI{57}: $\Tab_{AB}$.}
Here the outer pair is $\nabla\!\cdot\!\nabla'$ and $\Lambda_B=b=-k/(D-2)$:
\begin{equation*}
\Tab_{AB}^{\chIVloc}=\kappa^2\cdot i\,b\cdot\big(-a^{D-2}\nabla^2\big)\delta^D
=+\frac{i\kappa^2k}{D-2}\,a^{D-2}\nabla^2\delta^D
=+\frac{i\kappa^2H^2}{16\pi^2}\,a^2\nabla^2\delta^4(x-x')+\order{D-4}.
\tag{I.57$'$}\label{eq:ch04:57}
\end{equation*}

\paragraph{Eq.~\eqI{58}: $\Tab_{AC}$.}
With $\Lambda_C=c=k/[(D-2)(D-3)]$:
\begin{equation*}
\begin{aligned}
\Tab_{AC}^{\chIVloc}&=-2\kappa^2\,\frac{D-2}{D-3}\cdot i\,c\,\big(a^D\square-a^{D-2}\nabla^2\big)\delta^D
=-\frac{2i\kappa^2k}{(D-3)^2}\big(a^D\square-a^{D-2}\nabla^2\big)\delta^D\\
&=-\frac{i\kappa^2H^2}{4\pi^2}\big\{a^4\square-a^2\nabla^2\big\}\delta^4(x-x')+\order{D-4}.
\end{aligned}
\tag{I.58$'$}\label{eq:ch04:58}
\end{equation*}

In \eqref{eq:ch04:57} and \eqref{eq:ch04:58}, $b$ and $c$ are finite and are not multiplied by a pole.
Setting $D=4$ there therefore changes only terms of order $D-4$.
The pole coefficients in \eqref{eq:ch04:54} are kept at exact $D$.

% ---------------------------------------------------------------------
\subsection{Eq.~(I.59): the sum \texorpdfstring{\tD}{}}
% ---------------------------------------------------------------------

Adding \eqref{eq:ch04:54}, \eqref{eq:ch04:57} and \eqref{eq:ch04:58} gives, in units of $i\kappa^2k$ and at exact $D$,
\begin{equation}
\begin{aligned}
-iM^2_{\rm 3pt,\,loc}=i\kappa^2k\,\Big\{&\Big[\frac{D-1}{D-3}\,\pihat-\frac{2}{(D-3)^2}\Big]a^D\square\\
&+\Big[-\frac{D-1}{D-3}\,\pihat+\frac{1}{D-2}+\frac{2}{(D-3)^2}\Big]a^{D-2}\nabla^2\Big\}\delta^D(x-x').
\end{aligned}
\label{eq:ch04:59exact}
\end{equation}
Now we set $D=4$ in the finite parts: $\frac{2}{(D-3)^2}\to2$ and $\frac{1}{D-2}\to\frac12$.
The constant $k=\frac{H^{D-2}}{(4\pi)^{D/2}}\frac{\Gamma(D-1)}{\Gamma(D/2)}$ is the same prefactor as in the paper's \eqI{59}.
\begin{resultbox}[Result of Task 9: local 3-point contribution (definition sign)]
\begin{equation*}
\begin{aligned}
-iM^2_{\rm 3pt,\,loc}(x;x')=i\kappa^2k\,\Big\{&\Big[-\frac{D-1}{D-3}\,\pi\cot\big(\tfrac{\pi}{2}D\big)-2\Big]a^D\square\\
&+\Big[\frac{D-1}{D-3}\,\pi\cot\big(\tfrac{\pi}{2}D\big)+\frac52\Big]a^{D-2}\nabla^2+\order{D-4}\Big\}\delta^D(x-x').
\end{aligned}
\tag{I.59$'$}\label{eq:ch04:59}
\end{equation*}
\end{resultbox}
\noindent
Multiplying \eqref{eq:ch04:59} by $-1$ gives the paper's convention.
The brackets then read $\big[\frac{D-1}{D-3}\pi\cot(\frac\pi2D)+2\big]$ and $\big[-\frac{D-1}{D-3}\pi\cot(\frac\pi2D)-\frac52\big]$.
These differ from the printed \eqI{59} only in the finite constant of the second bracket: $-\frac52$ instead of $-\frac72$.

% ---------------------------------------------------------------------
\subsection{Comparison with Paper I}
\label{sec:ch04:t9comp}
% ---------------------------------------------------------------------

Derived expressions are in the definition sign; the paper's are as printed (paper's convention).

\begin{center}
\small
\renewcommand{\arraystretch}{1.35}
\begin{tabular}{@{}>{\raggedright\arraybackslash}p{2.7cm}>{\raggedright\arraybackslash}p{4.9cm}>{\raggedright\arraybackslash}p{4.9cm}>{\raggedright\arraybackslash}p{2.0cm}@{}}
\toprule
Item & Derived (this task) & Paper I (printed) & Status\\
\midrule
\eqI{54} & $\ldots\{+a^D\square-a^{D-2}\nabla^2\}\delta^D$ & $\ldots\{-a^D\square+a^{D-2}\nabla^2\}\delta^D$ & \chIVno\ sign\\
\eqI{55}--\eqI{56} & \eqref{eq:ch04:R2} & same & \chIVok\\
\eqI{57} & $+\frac{i\kappa^2H^2}{16\pi^2}\,a^2\nabla^2\delta^4$ & $-\frac{i\kappa^2H^2}{16\pi^2}\,a^2\nabla^2\delta^4$ & \chIVno\ sign\\
\eqI{58} & $-\frac{i\kappa^2H^2}{4\pi^2}\{a^4\square-a^2\nabla^2\}\delta^4$ & $+\frac{i\kappa^2H^2}{4\pi^2}\{a^4\square-a^2\nabla^2\}\delta^4$ & \chIVno\ sign\\
\eqI{59}, $a^D\square$ & $-\frac{D-1}{D-3}\pi\cot-2$ & $+\frac{D-1}{D-3}\pi\cot+2$ & \chIVno\ sign\\
\eqI{59}, $a^{D-2}\nabla^2$ & $+\frac{D-1}{D-3}\pi\cot+\frac52$ & $-\frac{D-1}{D-3}\pi\cot-\frac72$ & \chIVno\ sign and $\frac72\neq\frac52$\\
text: ``4th term of (47)'' & 2nd term & 4th term & \chIVno\ typo\\
\bottomrule
\end{tabular}
\end{center}

\paragraph{Diagnosis.}
\begin{enumerate}[leftmargin=1.5em,itemsep=2pt]
\item \textbf{Overall sign.} The left-hand sides of \eqI{54}, \eqI{57} and \eqI{58} carry the signs of the printed \eqI{47}--\eqI{49}.
The right-hand sides are correct evaluations of these left-hand sides.
The difference therefore comes entirely from the sign in the step \eqI{46}$\to$\eqI{47}--\eqI{49} (Task~7).
Consequently, the paper's 3-point contribution has the overall opposite sign relative to the definitions.
\item \textbf{Arithmetic of \eqI{59}.} The paper's own terms give $-\frac12$ [from \eqI{57}] $-\,2$ [from \eqI{58}] $=-\frac52$ for the finite part of the $a^{D-2}\nabla^2$ coefficient; the printed value is $-\frac72$.
In the paper's convention, the printed \eqI{59} minus the sum of the paper's own terms is $-\frac{i\kappa^2H^2}{8\pi^2}\,a^2\nabla^2\delta^4(x-x')$.
This term has the form of the counterterm \eqI{145}, whose vertex Paper~I prints as $i\alpha_3\kappa^2H^2a^{D-2}\nabla^2\delta^D(x-x')$.
It is absorbed by $\Delta\alpha_3=-\frac{1}{8\pi^2}$.
The paper treats the finite parts of the counterterms as arbitrary, so this error does not affect the physical result.
The paper most likely used this value in Section~IV, where it combines \eqI{44}, \eqI{59} and \eqI{67}.
\end{enumerate}

\paragraph{Effect on the accumulated findings.}
Task~6 concluded that ``in the paper, the 4-point contribution enters with the opposite sign relative to the 3-point contribution''.
That conclusion assumed that the paper's 3-point calculation used \eqI{46}.
The present task shows that the paper used \eqI{47}--\eqI{49} instead. Accordingly:
\begin{itemize}[leftmargin=1.2em,itemsep=2pt]
\item In the paper, the 4-point and 3-point contributions are \textbf{consistent} with each other.
Both have the opposite sign relative to the definitions \eqI{168} and \eqI{131}.
\item The paper's bare $M^2$ is $-1$ times the one defined by \eqI{168}.
The counterterms are chosen to cancel the divergences, so the renormalized result is also $-1$ times the correct one, up to arbitrary finite local terms.
This affects the sign of the mode-function correction in Ref.~[30] of Paper~I (= Paper~II, arXiv:0710.5282).
(Chapter~\ref{ch:sk} shows that at late times the non-falling total of that correction vanishes, so only decaying terms change sign and the late-time conclusion of Paper~II is unchanged.)
\item \textbf{Supporting hint.} After \eqI{67}, the paper states that each local term of \eqI{67} cancels a similar finite local 4-point term in \eqI{44}.
If the relative sign were inconsistent, this cancellation would not happen.
The cancellation also covers the $\ln(aa')$ and $9H^2a^D$ terms, which had raised a concern in Task~5 because they are not of counterterm form.
This remained to be confirmed once \eqI{60}--\eqI{67} were derived.
(It is confirmed in Task~10, Section~\ref{sec:ch04:cancel}.)
\end{itemize}

\medskip
\noindent\begin{minipage}{\linewidth}
\paragraph{Accumulated findings (Tasks 2--9).}\mbox{}
Identifiers in parentheses (S$n$, B$n$, C$n$) refer to the classified list of findings in Chapter~\ref{ch:summary}.
\begin{center}
\small
\renewcommand{\arraystretch}{1.3}
\begin{tabular}{@{}>{\raggedright\arraybackslash}p{2.4cm}>{\raggedright\arraybackslash}p{8.0cm}>{\raggedright\arraybackslash}p{4.0cm}@{}}
\toprule
Location & Finding & Effect\\
\midrule
\eqI{7}, left-hand side (C3) & $-\frac12\kappa$ printed; correct is $+\frac12\kappa$ & none\\
\eqI{14}, first coefficient (C4) & $\Gamma(\frac D2-1)/(\frac D2-1)$ printed; correct is $\Gamma(\frac D2-1)$ & none\\
Table II, $U_2$ (C5) & $-\frac i4$ printed; correct is $+\frac i4$ & none\\
\eqI{32}--\eqI{33} (S1) & overall sign of the 4-point contribution reversed & overall sign of the bare $M^2$\\
\eqI{39} (C6) & $\bar\eta_{\rho\sigma}\to\bar\eta^{\rho\sigma}$ & none\\
\eqI{40} (B1) & finite part $+2$ printed; correct is $-2$ & absorbable in counterterms\\
\eqI{47}--\eqI{49} (S1) & overall sign reversed in all nine terms & \textbf{carried into the later calculation} (this task)\\
\eqI{54}--\eqI{59} & overall sign reversed (inherited); in \eqI{59} $-\frac72\to-\frac52$; in the text ``4th''$\to$``2nd'' & sign: inherited (S1); \eqI{59}: absorbable in $\Delta\alpha_3$ (B2); ``4th'': typo (C7)\\
\textbf{Overall} (S1) & in the paper the 4-point and 3-point contributions are consistent with each other; both are reversed & bare $M^2=-$(definition)\\
\bottomrule
\end{tabular}
\end{center}
\end{minipage}


% ---------------------------------------------------------------------
\subsection{Code checks \texorpdfstring{\tC}{}}
\label{sec:ch04:t9code}
% ---------------------------------------------------------------------

The script \texttt{code/04\_three\_point\_local\_log/check\_54\_59.py} (SymPy; about 1\,s) checks the following.
It was actually run.
It is a check only and does not replace the analytic derivation above.
\begin{itemize}[leftmargin=1.2em,itemsep=3pt]
\item[(A)] \textbf{Local part directly from \eqI{45} and Table~I.}
This check does not use the intermediate step \eqI{47}--\eqI{49}.
The graviton tensor factors $T^A$, $T^B$, $T^C$ are built as explicit component arrays for $D=4,5,6$.
Outer derivatives at $x$ and $x'$ are represented by formal vectors $P$, $P'$, and inner derivatives by $Q$, $Q'$.
Each vertex factor is $i(-P^\alpha)Q^\beta-\frac i2\eta^{\alpha\beta}(-P\!\cdot\!Q)$, since partial integration turns an outer-leg derivative into $-P$.
The two vertex factors are contracted with $T^Ad_A+T^Bd_B+T^Cd_C$, where $d_A,d_B,d_C$ label the type of the outer propagator.
The inner derivatives are replaced by the delta-function part of \eqI{52}, i.e.\ $Q_\rho Q'_\sigma\to\delta^0_\rho\delta^0_\sigma$.
For $D=4,5,6$ the local polynomial comes out as
$\frac{D-1}{D-3}P_0P'_0\,d_A+(\bar P\!\cdot\!\bar P')\,d_B-\frac{2(D-2)}{D-3}P_0P'_0\,d_C$.
This is exactly the local part of the three terms in \eqref{eq:ch04:G1t9}~\chIVok.
\item[(B)] \textbf{Rules \eqref{eq:ch04:R2} and \eqI{55}--\eqI{56}.}
These are verified on a test function $g(\eta,x_1,x_2,x_3)$ with explicit $a(\eta)=-1/(H\eta)$ and symbolic $D$ in the powers of $a$~\chIVok.
\item[(C)] \textbf{The sum.}
The derived finite coefficients are $-2$ for $a^D\square$ and $+\frac52$ for $a^{D-2}\nabla^2$.
The paper's own terms add up to $-\frac{D-1}{D-3}\pihat+2$ for $a^D\square$ and $\frac{D-1}{D-3}\pihat-\frac52$ for $a^{D-2}\nabla^2$.
The printed $a^{D-2}\nabla^2$ value is $\frac{D-1}{D-3}\pihat-\frac72$.
The paper's \eqI{54}, \eqI{57} and \eqI{58} are $-1$ times the derived ones, term by term~\chIVok.
(For \eqI{54} this is checked at symbolic $D$; for \eqI{57} and \eqI{58}, which the paper prints only at $D=4$, the $D=4$ coefficients are compared.)
\end{itemize}

\paragraph{Actual output} (Python 3, SymPy 1.14.0; long lines wrapped, marked by \texttt{\textbackslash}).
The output of the re-run in the code folder is identical to the output in the original log.
{\scriptsize
\begin{verbatim}
== (A) local part extracted directly from (45) + Table I (no use of (47)-(49)) ==
  D=4: local polynomial == (D-1)/(D-3) P0P0' dA + (P.P')_spatial dB - 2(D-2)/(D-3) P0P0' dC : True
  D=5: local polynomial == (D-1)/(D-3) P0P0' dA + (P.P')_spatial dB - 2(D-2)/(D-3) P0P0' dC : True
  D=6: local polynomial == (D-1)/(D-3) P0P0' dA + (P.P')_spatial dB - 2(D-2)/(D-3) P0P0' dC : True
== (B) outer derivatives on a^{D-2} delta^D(x-x'), acting on a test function g (D=4 coordinates) ==
  d0 d0'[a^{D-2} delta] == (a^D Box - a^{D-2} Lap) delta : True
  sum_i d_i d_i'[a^{D-2} delta] == - a^{D-2} Lap delta   : True
  (55)-(56): -d0(a^{D-2} d0' delta) == -a^D Box + a^{D-2} Lap : True
== (C) assembling (59): units of i kappa^2 k, k = H^{D-2}/(4pi)^{D/2} Gamma(D-1)/Gamma(D/2) ==
  derived coeff of a^D Box     : (D**2*pihat - 4*D*pihat + 3*pihat - 2)/(D**2 - 6*D + 9)
  derived coeff of a^{D-2} Lap : (-D**3*pihat + 6*D**2*pihat + D**2 - 11*D*pihat - 4*D + 6*pihat + 5)\
                                 /(D**3 - 8*D**2 + 21*D - 18)
  finite parts at D=4 (pihat-terms kept exact):  Box: -2    Lap: 5/2
  sum of the paper's own (54)+(57)+(58):  Box: (-D*pihat + 2*D + pihat - 6)/(D - 3)    \
                                          Lap: (D*pihat - 5*D/2 - pihat + 15/2)/(D - 3)
  paper's printed (59):                    Box: -(D-1)/(D-3)*pihat + 2   Lap: (D-1)/(D-3)*pihat - 7/2
  paper's (54),(57),(58) = -(derived) term by term: True True True
  check k -> H^2/(8 pi^2) at D=4: H**2/(8*pi**2)
\end{verbatim}
}
\noindent
The two derived coefficients factor as $\frac{(D-1)(D-3)\pihat-2}{(D-3)^2}$ and
$\frac{-(D-1)(D-2)(D-3)\pihat+(D-3)^2+2(D-2)}{(D-2)(D-3)^2}$, which are the brackets of \eqref{eq:ch04:59exact}.
(In section (B) the label ``D=4 coordinates'' means three spatial coordinates; the exponent $D-2$ of $a$ is symbolic.)

% ---------------------------------------------------------------------
\subsection{Calculation record}
% ---------------------------------------------------------------------
\begin{itemize}[leftmargin=1.2em,itemsep=2pt]
\item \textbf{Derived.} Rules \eqref{eq:ch04:R1}--\eqref{eq:ch04:R2}; \eqI{55}--\eqI{56} confirmed;
\eqI{54} [Eq.~\eqref{eq:ch04:54}], \eqI{57} [Eq.~\eqref{eq:ch04:57}], \eqI{58} [Eq.~\eqref{eq:ch04:58}];
\eqI{59} at exact $D$ [Eq.~\eqref{eq:ch04:59exact}] and at $D=4$ in the finite parts [Eq.~\eqref{eq:ch04:59}].
\item \textbf{Checks passed.} SymPy (A): local part extracted directly from \eqI{45}; (B): outer-derivative rules; (C): the sum, and an audit of the paper's own arithmetic.
All were actually run.
\item \textbf{Discrepancies.} The overall sign is reversed. Its origin is \eqI{47}--\eqI{49}, whose printed signs the paper carried along.
In \eqI{59} the constant should be $-\frac52$ instead of $-\frac72$.
In the text, ``2nd term'' should replace ``4th term''.
\item \textbf{Open point (at the time of Task 9).} The logarithm contributions \eqI{60}--\eqI{67} and the expected cancellation with \eqI{44}.
Both are settled in Task~10.
\end{itemize}

% =====================================================================
\section{Task 10: the logarithm contributions, Eqs. (I.60)--(I.67)}
\label{sec:ch04:task10}
% =====================================================================

Task~10 derives the logarithm contributions of the 3-point diagram with the corrected sign and normalizations.
It then compares them with the paper (printed pages 124005-6 and 124005-7).

\begin{discbox}[Summary of findings (Task 10)]
\begin{itemize}[leftmargin=1.2em,itemsep=2pt]
\item \eqI{60}, \eqI{61} and \eqI{62} are reproduced exactly.
\eqI{62} was obtained with the corrected first coefficient of \eqI{14}; the printed coefficient gives the same result at $D=4$.
\item \eqI{63}--\eqI{66} were derived at general $D$. All four are exactly $-1$ times the paper's.
The paper's left-hand sides carry the signs of the printed \eqI{47}, and its right-hand sides correctly evaluate those left-hand sides.
This is the same diagnosis as in Task~9.
\item \eqI{67}: the derived result is $-1$ times the paper's.
The $\nabla^4A$ terms cancel exactly at $D=4$; at general $D$ the remaining coefficient is $-\frac{2(D-4)}{D-3}$.
\item \textbf{The cancellation against \eqI{44} is confirmed.}
In the correct \eqI{44}, the local terms containing $\ln(aa')$ and the $9H^2a^D$ term cancel against the local terms of the correct \eqI{67}.
At general $D$ they leave only a factor $\frac{(D-4)(D+1)}{4(D-3)}$.
This removes the concern left open in Task~5, which was that these terms are not of counterterm form.
It also confirms that the relative sign of the 4-point and 3-point contributions is the one fixed by the definitions.
In the paper both contributions are reversed, so the cancellation also happens in the paper.
\end{itemize}
\end{discbox}

% ---------------------------------------------------------------------
\subsection{Goal, inputs and scope}
% ---------------------------------------------------------------------

\paragraph{Target equations \tP.}
Paper~I states that all logarithm contributions come from \eqI{47} and can be simplified with the propagator equation \eqI{12}.
It prints (p.~124005-6, right column, and p.~124005-7, left column):
\begin{align*}
&\partial_0\big(a^{D-2}\partial_0 i\Delta_A(x;x')\big)=-i\delta^D(x-x')+a^{D-2}\nabla^2A(y),\tag{I.60}\\
&\partial'_0\big(a'^{D-2}\partial'_0 i\Delta_A(x;x')\big)=-i\delta^D(x-x')+a'^{D-2}\nabla^2A(y),\tag{I.61}\\
&A(y)=\frac{H^2}{16\pi^2}\Big\{\frac4y-2\ln\Big(\frac y4\Big)-1+\order{D-4}\Big\},\tag{I.62}\\[3pt]
&\kappa^2\nabla\!\cdot\!\nabla'\big[(aa')^{D-2}k\ln(aa')\nabla\!\cdot\!\nabla' i\Delta_A(x;x')\big]
=\frac{\kappa^2H^2}{8\pi^2}\ln(aa')(aa')^2\nabla^4A(y)+\order{D-4},\tag{I.63}\\[3pt]
&-\kappa^2\Big(\frac{D-1}{D-3}\Big)\partial_0\partial'_0\big[(aa')^{D-2}k\ln(aa')\partial_0\partial'_0 i\Delta_A(x;x')\big]\notag\\
&\qquad=\frac{i3\kappa^2H^2}{8\pi^2}\ln(aa')\{-a^4\square+2a^2\nabla^2\}\delta^4(x-x')\notag\\
&\qquad-\frac{3\kappa^2H^2}{8\pi^2}\ln(aa')(aa')^2\nabla^4A(y)-\frac{i9\kappa^2H^4}{8\pi^2}a^4\delta^4(x-x')\notag\\
&\qquad-\frac{3\kappa^2H^3}{8\pi^2}(aa')^2(a\partial_0+a'\partial'_0)\nabla^2A(y)+\order{D-4},\tag{I.64}\\[3pt]
&\kappa^2\partial_i\partial'_0\big[(aa')^{D-2}k\ln(aa')\partial_i\partial'_0 i\Delta_A(x;x')\big]
=-\frac{i\kappa^2H^2}{8\pi^2}\ln(aa')a^2\nabla^2\delta^4(x-x')\notag\\
&\qquad+\frac{\kappa^2H^2}{8\pi^2}\ln(aa')(aa')^2\nabla^4A(y)+\frac{\kappa^2H^3}{8\pi^2}(aa')^2a'\partial'_0\nabla^2A(y)+\order{D-4},\tag{I.65}\\[3pt]
&\kappa^2\partial_0\partial'_i\big[(aa')^{D-2}k\ln(aa')\partial_0\partial'_i i\Delta_A(x;x')\big]
=-\frac{i\kappa^2H^2}{8\pi^2}\ln(aa')a^2\nabla^2\delta^4(x-x')\notag\\
&\qquad+\frac{\kappa^2H^2}{8\pi^2}\ln(aa')(aa')^2\nabla^4A(y)+\frac{\kappa^2H^3}{8\pi^2}(aa')^2a\partial_0\nabla^2A(y)+\order{D-4},\tag{I.66}\\[3pt]
&-iM^2_{\rm 3pt,\,log}(x;x')=\frac{\kappa^2H^2}{8\pi^2}\Big\{\ln(aa')\big[-3a^4\square+4a^2\nabla^2\big]i\delta^4(x-x')\notag\\
&\qquad-9H^2a^4i\delta^4(x-x')-2H(aa')^2(a\partial_0+a'\partial'_0)\nabla^2A(y)+\order{D-4}\Big\}.\tag{I.67}
\end{align*}
Between \eqI{61} and \eqI{62} the paper states that the limit $D=4$ may be taken because all logarithm contributions are finite.
None of these printed expressions was used as an input. They enter only the comparison in Section~\ref{sec:ch04:t10comp}.

\paragraph{Inputs.}
\begin{itemize}[leftmargin=1.2em,itemsep=3pt]
\item[(G1)] The derived \eqI{47} \tD\ (Task~7, the $AA$ terms).
In the logarithm contribution, the outer $i\Delta_A$ is replaced by $k\ln(aa')$.
Abbreviation: $\mathcal P\equiv(aa')^{D-2}\,k\,L$. Then
\begin{equation}
\begin{aligned}
\mathcal L_1&=-\kappa^2\,\nabla\!\cdot\!\nabla'\big[\mathcal P\,\nabla\!\cdot\!\nabla' i\Delta_A\big],&
\mathcal L_2&=+\kappa^2\,\frac{D-1}{D-3}\,\partial_0\partial'_0\big[\mathcal P\,\partial_0\partial'_0 i\Delta_A\big],\\
\mathcal L_3&=-\kappa^2\,\partial_i\partial'_0\big[\mathcal P\,\partial_i\partial'_0 i\Delta_A\big],&
\mathcal L_4&=-\kappa^2\,\partial_0\partial'_i\big[\mathcal P\,\partial_0\partial'_i i\Delta_A\big].
\end{aligned}
\label{eq:ch04:G1t10}
\end{equation}
(Repeated spatial indices $i$ are summed.)
In the paper's \eqI{47} the signs are $+,-,+,+$ in this order, i.e.\ reversed.
\item[(G2)] The propagator equation \eqI{12} \tP, together with \eqI{13} and \eqI{14} \tP.
Following Task~8, the first coefficient of \eqI{14} is taken to be $\Gamma(\frac D2-1)$.
\item[(G3)] Homogeneity \tD\ (Task~7): $\partial'_i\,i\Delta_A=-\partial_i\,i\Delta_A$.
Hence $\nabla\!\cdot\!\nabla' i\Delta_A=-\nabla^2A$, and $\nabla'^2=\nabla^2$ on functions of $\vec x-\vec x'$.
\item[(G4)] Separation convention \tP\ (Task~9, (G4)).
The logarithm contribution is $k\ln(aa')$ times the \emph{full} $\partial\partial' i\Delta_A$, i.e.\ including the local term of \eqI{52}.
Together with the local contribution of Task~9, the decomposition is then complete.
\end{itemize}

\paragraph{Scope.}
Not treated here: the normal contributions (\eqI{68} and later), and the elimination of time derivatives that the paper performs after \eqI{67}.
The calculation is done at general $D$; the passage to $D=4$ is justified in Section~\ref{sec:ch04:Dto4}.

% ---------------------------------------------------------------------
\subsection{Eqs.~(I.60) and (I.61) \texorpdfstring{\tD}{}}
% ---------------------------------------------------------------------

By \eqI{9}, $\sqrt{-\hat g}\,\hat g^{\mu\nu}=a^{D-2}\eta^{\mu\nu}$, so the propagator equation \eqI{12} reads
$-\partial_0\big(a^{D-2}\partial_0 i\Delta_A\big)+a^{D-2}\nabla^2 i\Delta_A=i\delta^D(x-x')$.
Since $\nabla^2[k\ln(aa')]=0$, we have $\nabla^2 i\Delta_A=\nabla^2A$. Hence
\begin{equation*}
\boxed{\;\partial_0\big(a^{D-2}\partial_0\,i\Delta_A(x;x')\big)=-i\delta^D(x-x')+a^{D-2}\nabla^2A(y)\;}
\tag{I.60}\label{eq:ch04:60}
\end{equation*}
Because of the symmetry $i\Delta_A(x;x')=i\Delta_A(x';x)$, the same equation holds at $x'$. This is \eqI{61}:
\begin{equation*}
\partial'_0\big(a'^{D-2}\partial'_0\,i\Delta_A\big)=-i\delta^D(x-x')+a'^{D-2}\nabla^2A\quad\chIVok
\tag{I.61}\label{eq:ch04:61}
\end{equation*}

\paragraph{Corollary.}
We have $\partial_0\big(a^p\partial_0 l\big)=\partial_0\big(Ha^{D-1}\big)=(D-1)H^2a^D$.
With $A=i\Delta_A-kL$ and $\partial_0L=\partial_0l$, this gives
\begin{equation}
\partial_0\big(a^p\partial_0A\big)=-i\delta^D+a^p\nabla^2A-(D-1)kH^2a^D .
\end{equation}
Applying $\nabla^2$ removes the last term:
\begin{equation}
\partial_0\big(a^p\partial_0\nabla^2A\big)=-i\nabla^2\delta^D+a^p\nabla^4A .
\label{eq:ch04:corol}
\end{equation}
The primed version, $\partial'_0\big(a'^p\partial'_0\nabla^2A\big)=-i\nabla^2\delta^D+a'^p\nabla^4A$, follows in the same way from \eqref{eq:ch04:61}.
Both forms are used repeatedly below.

% ---------------------------------------------------------------------
\subsection{Eq.~(I.62): \texorpdfstring{$A(y)$}{A(y)} at \texorpdfstring{$D=4$}{D=4} \texorpdfstring{\tD}{}}
% ---------------------------------------------------------------------

Write $D=4+\epsilon$ and $\mathcal N\equiv\frac{H^{D-2}}{(4\pi)^{D/2}}\to\frac{H^2}{16\pi^2}$.
We expand the terms of \eqI{14} one at a time:
\begin{equation}
\begin{aligned}
&\Gamma\big(\tfrac D2-1\big)\Big(\frac4y\Big)^{\frac D2-1}\to\frac4y,\\
&\frac{\Gamma(\frac D2+1)}{\frac D2-2}\Big(\frac4y\Big)^{\frac D2-2}=\frac4\epsilon+2\psi(3)+2\ln\frac4y+\order{\epsilon}
=\frac4\epsilon+3-2\gamma+2\ln\frac4y+\order{\epsilon},\\
&-\pi\cot\big(\tfrac{\pi D}{2}\big)\frac{\Gamma(D-1)}{\Gamma(\frac D2)}=\Big(-\frac2\epsilon+\order{\epsilon}\Big)\cdot2\Big[1+\epsilon\Big(1-\frac\gamma2\Big)\Big]
=-\frac4\epsilon-4+2\gamma+\order{\epsilon}.
\end{aligned}
\end{equation}
Here we used $\psi(3)=\frac32-\gamma$ and $\psi(2)=1-\gamma$; the bracket in the last line is
$\Gamma(3+\epsilon)/\Gamma(2+\frac\epsilon2)=2[1+\epsilon(\psi(3)-\frac12\psi(2))]$.
The $n$-th term of the series in \eqI{14} vanishes identically at $D=4$: both of its parts equal $\frac{n+2}{n}(y/4)^n$ there.
The poles cancel exactly, with the same factor $\mathcal N$ in front of both. Therefore
\begin{equation*}
\boxed{\;A(y)=\frac{H^2}{16\pi^2}\Big\{\frac4y-2\ln\frac y4-1+\order{D-4}\Big\}\;}
\tag{I.62}\label{eq:ch04:62}
\end{equation*}
This is the paper's \eqI{62}~\chIVok.
The printed first coefficient $\Gamma(\frac D2-1)/(\frac D2-1)$ also tends to 1 at $D=4$, so it affects \eqI{62} only at order $D-4$.

\emph{Check.} $(4y-y^2)A''+4(2-y)A'=6h=3k$ with $h\equiv H^2/(16\pi^2)$.
This is $\square A=(D-1)kH^2$ at $D=4$, as expected from $\square\,i\Delta_A=0$ away from coincidence and $\square\,k\ln(aa')=-(D-1)kH^2$.

% ---------------------------------------------------------------------
\subsection{Distributional tools \texorpdfstring{\tD}{}}
% ---------------------------------------------------------------------

\begin{itemize}[leftmargin=1.5em,itemsep=4pt]
\item[(D1)] For a function $f$ of time only (repeating Lemma~B of Task~5):
\begin{equation}
f(\eta')\,\partial_0\delta^D=f(\eta)\,\partial_0\delta^D+f'(\eta)\,\delta^D,\qquad
f(\eta')\,\partial_0^2\delta^D=f\,\partial_0^2\delta^D+2f'\,\partial_0\delta^D+f''\,\delta^D .
\label{eq:ch04:D1}
\end{equation}
(Acting on a test function $h(x')$, the left-hand side of the second identity gives $\partial_0^2(fh)=fh''+2f'h'+f''h$.)
\item[(D2)] Symmetry of the kernel:
\begin{equation}
\partial'_0\big[F(\eta')\,\partial'_0\delta^D(x-x')\big]\equiv\partial_0\big[F(\eta)\,\partial_0\delta^D(x-x')\big].
\label{eq:ch04:D2}
\end{equation}
Both act on a test function $g(x')$ as $\partial_0(F\partial_0g)$ (integrate by parts twice).
\item[(D3)] $f(\eta,\eta')\,\nabla^2\delta^D(x-x')=f(\eta,\eta)\,\nabla^2\delta^D$, because only spatial derivatives act on the delta function.
For example, $a'^p\,l\,\nabla^2\delta^D=a^p\,l\,\nabla^2\delta^D$ and $L\,\nabla^2\delta^D=2l\,\nabla^2\delta^D$.
\item[(D4)] $\partial_0\big(a^p\partial_0\delta^D\big)=-\big(a^D\square-a^{D-2}\nabla^2\big)\delta^D$, from \eqI{42} (Task~9, Eq.~\eqref{eq:ch04:R2}).
\item[(D5)] \textbf{Lemma.}
\begin{equation}
-2i\,\partial_0\big[a^p l\,\partial_0\delta^D\big]=i\,L\,\big(a^D\square-a^{D-2}\nabla^2\big)\delta^D+i(D-1)H^2a^D\delta^D .
\label{eq:ch04:D5}
\end{equation}
\emph{Proof.} First, $\partial_0[a^pl\,\partial_0\delta]=l\,\partial_0(a^p\partial_0\delta)+Ha^{D-1}\partial_0\delta$.
Next apply (D1) with $f=l(\eta')$, for which $f'=Ha$ and $f''=H^2a^2$, together with $\partial_0a^p=pHa^{D-1}$:
\begin{equation}
\begin{aligned}
l'\,\partial_0(a^p\partial_0\delta)&=l'\big(pHa^{D-1}\partial_0\delta+a^p\partial_0^2\delta\big)\\
&=pHa^{D-1}\big(l\,\partial_0\delta+Ha\,\delta\big)+a^p\big(l\,\partial_0^2\delta+2Ha\,\partial_0\delta+H^2a^2\delta\big)\\
&=l\,\partial_0(a^p\partial_0\delta)+2Ha^{D-1}\partial_0\delta+(D-1)H^2a^D\delta .
\end{aligned}
\end{equation}
Therefore $L\,\partial_0(a^p\partial_0\delta)=2\,\partial_0[a^pl\,\partial_0\delta]+(D-1)H^2a^D\delta$.
Applying (D4) to the left-hand side gives \eqref{eq:ch04:D5}.~$\blacksquare$
\end{itemize}

% ---------------------------------------------------------------------
\subsection{The four terms at general \texorpdfstring{$D$}{D} \texorpdfstring{\tD}{}}
% ---------------------------------------------------------------------

\paragraph{Eq.~\eqI{63}: $\mathcal V_1\equiv\nabla\!\cdot\!\nabla'\big[\mathcal P\,\nabla\!\cdot\!\nabla' i\Delta_A\big]$.}
By (G3), $\nabla\!\cdot\!\nabla' i\Delta_A=-\nabla^2A$.
$\mathcal P$ is constant in space, and $\nabla\!\cdot\!\nabla'F(\vec x-\vec x')=-\nabla^2F$. Hence
\begin{equation*}
\mathcal V_1=k\,(aa')^{D-2}L\,\nabla^4A,\qquad
\mathcal L_1=-\kappa^2\mathcal V_1\;\xrightarrow{D=4}\;-\frac{\kappa^2H^2}{8\pi^2}\ln(aa')(aa')^2\nabla^4A(y).
\tag{I.63$'$}\label{eq:ch04:63}
\end{equation*}
This term has no local part, because it contains no time derivatives.
(In the other three terms the delta functions arise through \eqref{eq:ch04:60}--\eqref{eq:ch04:61} and the corollary \eqref{eq:ch04:corol}, i.e.\ from two time derivatives acting at the same point.)

\paragraph{Eq.~\eqI{64}: $\mathcal W\equiv\partial_0\partial'_0\big[\mathcal P\,\partial_0\partial'_0 i\Delta_A\big]$.}
Split $L=l+l'$.
In the $l$ part, $a^pl$ depends only on $\eta$, so $\partial'_0$ passes through it.
Moreover $a'^p\,\partial_0\partial'_0 i\Delta_A=\partial_0\big(a'^p\partial'_0 i\Delta_A\big)$, because $a'$ does not depend on $\eta$.
Using the commutativity of the derivatives and \eqref{eq:ch04:61}:
\begin{equation}
\partial_0\partial'_0\big[a^pl\cdot a'^p\,\partial_0\partial'_0 i\Delta_A\big]
=\partial_0\big[a^pl\;\partial_0\,\partial'_0\big(a'^p\partial'_0 i\Delta_A\big)\big]
=\partial_0\big[a^pl\;\partial_0\big(-i\delta^D+a'^p\nabla^2A\big)\big].
\end{equation}
The $l'$ part is the mirror image under $x\leftrightarrow x'$: $\partial'_0\big[a'^pl'\,\partial'_0\big(-i\delta^D+a^p\nabla^2A\big)\big]$.
\begin{itemize}[leftmargin=1.2em,itemsep=3pt]
\item \textbf{Local part.} By (D2) and (D5):
\begin{equation}
\begin{aligned}
-i\big\{\partial_0[a^pl\,\partial_0\delta]+\partial'_0[a'^pl'\,\partial'_0\delta]\big\}&=-2i\,\partial_0[a^pl\,\partial_0\delta]\\
&=iL\big(a^D\square-a^{D-2}\nabla^2\big)\delta^D+i(D-1)H^2a^D\delta^D .
\end{aligned}
\end{equation}
\item \textbf{$A$ part.} By the corollary \eqref{eq:ch04:corol},
$\partial_0[a^pl\,\partial_0\nabla^2A]=l\big(-i\nabla^2\delta+a^p\nabla^4A\big)+Ha^{D-1}\partial_0\nabla^2A$.
The mirror term is analogous. With (D3) for the delta terms, and $a'^pa^{D-1}=(aa')^pa$:
\begin{equation}
\begin{aligned}
&a'^p\,\partial_0[a^pl\,\partial_0\nabla^2A]+a^p\,\partial'_0[a'^pl'\,\partial'_0\nabla^2A]\\
&\qquad=-ia^{D-2}L\,\nabla^2\delta^D+(aa')^{D-2}L\,\nabla^4A+H(aa')^{D-2}\big(a\partial_0+a'\partial'_0\big)\nabla^2A .
\end{aligned}
\end{equation}
\end{itemize}
Adding both parts:
\begin{equation}
\begin{aligned}
\frac{\mathcal W}{k}={}&iL\big(a^D\square-2a^{D-2}\nabla^2\big)\delta^D+i(D-1)H^2a^D\delta^D\\
&+(aa')^{D-2}L\,\nabla^4A+H(aa')^{D-2}\big(a\partial_0+a'\partial'_0\big)\nabla^2A .
\end{aligned}
\label{eq:ch04:W}
\end{equation}
At $D=4$, $\mathcal L_2=+3\kappa^2\mathcal W$ with $k=\frac{H^2}{8\pi^2}$:
\begin{equation*}
\begin{aligned}
\mathcal L_2={}&\frac{3i\kappa^2H^2}{8\pi^2}\ln(aa')\{a^4\square-2a^2\nabla^2\}\delta^4+\frac{3\kappa^2H^2}{8\pi^2}\ln(aa')(aa')^2\nabla^4A\\
&+\frac{9i\kappa^2H^4}{8\pi^2}a^4\delta^4+\frac{3\kappa^2H^3}{8\pi^2}(aa')^2\big(a\partial_0+a'\partial'_0\big)\nabla^2A .
\end{aligned}
\tag{I.64$'$}\label{eq:ch04:64}
\end{equation*}

\paragraph{Eq.~\eqI{65}: $\mathcal V_3\equiv\partial_i\partial'_0\big[\mathcal P\,\partial_i\partial'_0 i\Delta_A\big]$.}
The outer $\partial_i$ passes through $\mathcal P$, which is constant in space, and meets the inner $\partial_i$ to form $\nabla^2$:
\begin{equation}
\mathcal V_3=k\,\partial'_0\big[(aa')^{D-2}L\,\partial'_0\nabla^2A\big]
=k\,a^p\Big\{\partial'_0\big[a'^pl'\,\partial'_0\nabla^2A\big]+l\,\partial'_0\big[a'^p\partial'_0\nabla^2A\big]\Big\}.
\end{equation}
Applying $\nabla^2$ to \eqref{eq:ch04:61} gives $\partial'_0\big(a'^p\partial'_0\nabla^2A\big)=-i\nabla^2\delta+a'^p\nabla^4A$.
Also, by the product rule with $\partial'_0l'=Ha'$,
\begin{equation}
\partial'_0\big[a'^pl'\,\partial'_0\nabla^2A\big]=l'\,\partial'_0\big(a'^p\partial'_0\nabla^2A\big)+Ha'^{D-1}\partial'_0\nabla^2A .
\end{equation}
Hence
\begin{equation}
\frac{\mathcal V_3}{k}=-ia^{D-2}L\,\nabla^2\delta^D+(aa')^{D-2}L\,\nabla^4A+H(aa')^{D-2}a'\partial'_0\nabla^2A .
\label{eq:ch04:V3}
\end{equation}

\paragraph{Eq.~\eqI{66}: $\mathcal V_4\equiv\partial_0\partial'_i\big[\mathcal P\,\partial_0\partial'_i i\Delta_A\big]$.}
Using $\nabla'^2=\nabla^2$, we obtain the mirror image of \eqref{eq:ch04:V3}:
\begin{equation}
\frac{\mathcal V_4}{k}=-ia^{D-2}L\,\nabla^2\delta^D+(aa')^{D-2}L\,\nabla^4A+H(aa')^{D-2}a\,\partial_0\nabla^2A .
\label{eq:ch04:V4}
\end{equation}
With the derived signs, $\mathcal L_3=-\kappa^2\mathcal V_3$ and $\mathcal L_4=-\kappa^2\mathcal V_4$. At $D=4$:
\begin{equation*}
\begin{aligned}
\mathcal L_3={}&+\frac{i\kappa^2H^2}{8\pi^2}\ln(aa')a^2\nabla^2\delta^4-\frac{\kappa^2H^2}{8\pi^2}\ln(aa')(aa')^2\nabla^4A\\
&-\frac{\kappa^2H^3}{8\pi^2}(aa')^2a'\partial'_0\nabla^2A,
\end{aligned}
\tag{I.65$'$}\label{eq:ch04:65}
\end{equation*}
and $\mathcal L_4$ is the same except that $a'\partial'_0$ is replaced by $a\partial_0$ in the last term:
\begin{equation*}
\begin{aligned}
\mathcal L_4={}&+\frac{i\kappa^2H^2}{8\pi^2}\ln(aa')a^2\nabla^2\delta^4-\frac{\kappa^2H^2}{8\pi^2}\ln(aa')(aa')^2\nabla^4A\\
&-\frac{\kappa^2H^3}{8\pi^2}(aa')^2a\,\partial_0\nabla^2A .
\end{aligned}
\tag{I.66$'$}\label{eq:ch04:66}
\end{equation*}

% ---------------------------------------------------------------------
\subsection{The sum at general \texorpdfstring{$D$}{D} \texorpdfstring{\tD}{}}
% ---------------------------------------------------------------------

The logarithm contribution is
\begin{equation}
-iM^2_{\rm 3pt,\,log}=\kappa^2\Big[-\mathcal V_1+\frac{D-1}{D-3}\,\mathcal W-\mathcal V_3-\mathcal V_4\Big].
\end{equation}
The following table collects the coefficient of each structure, in units of $\kappa^2k$, from \eqref{eq:ch04:63}, \eqref{eq:ch04:W}, \eqref{eq:ch04:V3} and \eqref{eq:ch04:V4}.
\begin{center}
\small
\renewcommand{\arraystretch}{1.45}
\begin{tabular}{@{}lcccccc@{}}
\toprule
Structure $/\kappa^2k$ & $-\mathcal V_1$ & $\frac{D-1}{D-3}\mathcal W$ & $-\mathcal V_3$ & $-\mathcal V_4$ & Total & $D=4$\\
\midrule
$L\,a^D\square\,i\delta^D$ & --- & $\frac{D-1}{D-3}$ & --- & --- & $\frac{D-1}{D-3}$ & $3$\\
$L\,a^{D-2}\nabla^2\,i\delta^D$ & --- & $-\frac{2(D-1)}{D-3}$ & $+1$ & $+1$ & $-\frac{4}{D-3}$ & $-4$\\
$H^2a^D\,i\delta^D$ & --- & $\frac{(D-1)^2}{D-3}$ & --- & --- & $\frac{(D-1)^2}{D-3}$ & $9$\\
$(aa')^{D-2}L\,\nabla^4A$ & $-1$ & $\frac{D-1}{D-3}$ & $-1$ & $-1$ & $-\frac{2(D-4)}{D-3}$ & $0$\\
$H(aa')^{D-2}a\partial_0\nabla^2A$ (+ mirror) & --- & $\frac{D-1}{D-3}$ & $-1$ & $-1$ & $\frac{2}{D-3}$ & $2$\\
\bottomrule
\end{tabular}
\end{center}
In the last row, $-\mathcal V_4$ contributes to the $a\partial_0$ term and $-\mathcal V_3$ to its mirror, the $a'\partial'_0$ term.
So each of the two terms receives $\frac{D-1}{D-3}-1=\frac{2}{D-3}$.
The arithmetic of the other totals is:
$-\frac{2(D-1)}{D-3}+2=-\frac{4}{D-3}$ and $\frac{D-1}{D-3}-3=\frac{8-2D}{D-3}=-\frac{2(D-4)}{D-3}$.
Altogether:
\begin{equation}
\begin{aligned}
-iM^2_{\rm 3pt,\,log}=\kappa^2k\,\Big\{&\ln(aa')\Big[\frac{D-1}{D-3}\,a^D\square-\frac{4}{D-3}\,a^{D-2}\nabla^2\Big]i\delta^D+\frac{(D-1)^2}{D-3}\,H^2a^D\,i\delta^D\\
&-\frac{2(D-4)}{D-3}\,(aa')^{D-2}\ln(aa')\,\nabla^4A+\frac{2}{D-3}\,H(aa')^{D-2}\big(a\partial_0+a'\partial'_0\big)\nabla^2A\Big\}.
\end{aligned}
\label{eq:ch04:67D}
\end{equation}

% ---------------------------------------------------------------------
\subsection{Passage to \texorpdfstring{$D=4$}{D=4} and its justification \texorpdfstring{\tA}{}}
\label{sec:ch04:Dto4}
% ---------------------------------------------------------------------

Every term in \eqref{eq:ch04:67D} is of one of two kinds:
\begin{enumerate}[leftmargin=1.5em,itemsep=1pt]
\item a $\delta^D$-local term whose coefficient is finite and has no pole;
\item a smooth function (such as $(aa')^{D-2}\ln(aa')$) times derivatives of a \emph{single} propagator.
\end{enumerate}
$A(y)\sim y^{1-D/2}\sim\Delta x^{2-D}$ is locally integrable in $D$ dimensions.
Its derivatives are therefore well-defined distributions; they are analytic near $D=4$ and carry no pole.
This is not the case for \emph{products of two propagators}, such as the $y^{2-D}$ that occurs in the normal contributions.
Consequently, the $\nabla^4A$ term with coefficient of order $D-4$ vanishes as $D\to4$.
Setting $D=4$ in the other coefficients changes them only at order $D-4$.
This confirms the paper's statement that all logarithm contributions are finite. The result is:
\begin{resultbox}[Result of Task 10: logarithm 3-point contribution (definition sign)]
\begin{equation*}
\begin{aligned}
-iM^2_{\rm 3pt,\,log}(x;x')=\frac{\kappa^2H^2}{8\pi^2}\Big\{&\ln(aa')\big[3a^4\square-4a^2\nabla^2\big]i\delta^4(x-x')+9H^2a^4\,i\delta^4(x-x')\\
&+2H(aa')^2\big(a\partial_0+a'\partial'_0\big)\nabla^2A(y)+\order{D-4}\Big\}
\end{aligned}
\tag{I.67$'$}\label{eq:ch04:67}
\end{equation*}
\end{resultbox}

% ---------------------------------------------------------------------
\subsection{Cancellation against (I.44) at general \texorpdfstring{$D$}{D} \texorpdfstring{\tD}{}}
\label{sec:ch04:cancel}
% ---------------------------------------------------------------------

The full-$D$ form of the correct \eqI{44} was derived in Task~5. It is
\begin{equation}
\begin{aligned}
-iM^2_{\rm 4pt}=-i\kappa^2k\Big\{&\tfrac14D(D-1)\big[\pihat+\ln(aa')\big]a^D\square-D\big[\pihat+\ln(aa')\big]a^{D-2}\nabla^2\\
&+\tfrac14D(D-1)^2H^2a^D\Big\}\delta^D+\order{D-4}.
\end{aligned}
\end{equation}
Its local terms containing $\ln(aa')$, together with its $H^2$ term, are
\begin{equation}
-iM^2_{\rm 4pt}\supset-i\kappa^2k\Big\{\tfrac14D(D-1)\ln(aa')\,a^D\square-D\ln(aa')\,a^{D-2}\nabla^2+\tfrac14D(D-1)^2H^2a^D\Big\}\delta^D .
\end{equation}
Adding the local terms of \eqref{eq:ch04:67D} gives
\begin{equation}
-i\kappa^2k\,\frac{(D-4)(D+1)}{4(D-3)}\Big\{(D-1)\ln(aa')\,a^D\square-4\ln(aa')\,a^{D-2}\nabla^2+(D-1)^2H^2a^D\Big\}\delta^D\;\xrightarrow{D\to4}\;0 .
\label{eq:ch04:cancel}
\end{equation}
The identities used are
\begin{equation}
\frac{1}{D-3}-\frac D4=-\frac{(D-4)(D+1)}{4(D-3)},\qquad D-\frac{4}{D-3}=\frac{(D-4)(D+1)}{D-3}.
\end{equation}
(For example, the $\ln(aa')\,a^D\square$ coefficient is $(D-1)\big[\frac{1}{D-3}-\frac D4\big]$ in units of $i\kappa^2k$, and $D^2-3D-4=(D-4)(D+1)$.)
The expression \eqref{eq:ch04:cancel} is not multiplied by a pole, so it vanishes as $D\to4$. Thus:
\begin{itemize}[leftmargin=1.2em,itemsep=2pt]
\item The terms $3\ln(aa')a^D\square$, $-4\ln(aa')a^{D-2}\nabla^2$ and $9H^2a^D$ cancel exactly.
In Task~5 these terms had raised a concern because they are not of counterterm form.
\item The cancellation confirms that the relative sign of the 4-point and 3-point contributions is the one fixed by the definitions.
With the other relative sign, these terms would add instead of cancelling.
\item In the paper both contributions are reversed.
So the paper's statement after \eqI{67}, that each local term of \eqI{67} cancels a similar finite local 4-point term in \eqI{44}, is also correct in the paper's own convention.
\item The only local 4-point remainder that is not cancelled consists of the $\mp2$ terms coming from the error in \eqI{40} (Task~5).
These are of counterterm form.
(They are present in the paper's \eqI{44} and absent from the derived one; see Task~5. The statement concerns the finite local terms: the divergent terms $\propto\pihat$ of \eqI{44} do not take part in this cancellation.)
\end{itemize}

% ---------------------------------------------------------------------
\subsection{Comparison with Paper I}
\label{sec:ch04:t10comp}
% ---------------------------------------------------------------------

Derived expressions are in the definition sign; the paper's are as printed (paper's convention).

\begin{center}
\small
\renewcommand{\arraystretch}{1.35}
\begin{tabular}{@{}>{\raggedright\arraybackslash}p{2.6cm}>{\raggedright\arraybackslash}p{5.4cm}>{\raggedright\arraybackslash}p{4.6cm}>{\raggedright\arraybackslash}p{1.8cm}@{}}
\toprule
Item & Derived (this task) & Paper I (printed) & Status\\
\midrule
\eqI{60}, \eqI{61} & \eqref{eq:ch04:60}, \eqref{eq:ch04:61} & same & \chIVok\\
\eqI{62} & \eqref{eq:ch04:62} & same & \chIVok\\
\eqI{63} & $-\frac{\kappa^2H^2}{8\pi^2}L(aa')^2\nabla^4A$ & $+\frac{\kappa^2H^2}{8\pi^2}L(aa')^2\nabla^4A$ & \chIVno\ sign\\
\eqI{64} & \eqref{eq:ch04:64} & all four terms reversed & \chIVno\ sign\\
\eqI{65} & $-\kappa^2\mathcal V_3$ & $+\kappa^2\mathcal V_3$ & \chIVno\ sign\\
\eqI{66} & $-\kappa^2\mathcal V_4$ & $+\kappa^2\mathcal V_4$ & \chIVno\ sign\\
\eqI{67} & $\{L[3a^4\square-4a^2\nabla^2]i\delta+9H^2a^4i\delta$\newline\hspace*{1em}$+2H(\ldots)\nabla^2A\}$ & $\{L[-3a^4\square+4a^2\nabla^2]i\delta$\newline\hspace*{1em}$-9H^2a^4i\delta-2H(\ldots)\nabla^2A\}$ & \chIVno\ sign\\
cancellation of $\nabla^4A$; coefficients 3, 4, 9, 2 & same & same & \chIVok\\
local cancellation with \eqI{44} & yes, \eqref{eq:ch04:cancel} & yes (text) & \chIVok\\
\bottomrule
\end{tabular}
\end{center}

\begin{notebox}[Diagnosis]
The left-hand sides of the paper's \eqI{63}--\eqI{66} carry the signs of the printed \eqI{47}, and the right-hand sides are correct evaluations of them.
The paper's \eqI{67} is the correct sum of its own \eqI{63}--\eqI{66}.
The difference comes entirely from the sign in the step \eqI{46}$\to$\eqI{47} (Task~7).
No new error was found in this task.
\end{notebox}

\medskip
\noindent\begin{minipage}{\linewidth}
\paragraph{Accumulated findings (Tasks 2--10).}\mbox{}
\begin{center}
\small
\renewcommand{\arraystretch}{1.3}
\begin{tabular}{@{}>{\raggedright\arraybackslash}p{2.4cm}>{\raggedright\arraybackslash}p{8.0cm}>{\raggedright\arraybackslash}p{4.0cm}@{}}
\toprule
Location & Finding & Effect\\
\midrule
\eqI{7}, left-hand side (C3) & $-\frac12\kappa$ printed; correct is $+\frac12\kappa$ & none\\
\eqI{14}, first coefficient (C4) & $\Gamma(\frac D2-1)/(\frac D2-1)$ printed; correct is $\Gamma(\frac D2-1)$ & none\\
Table II, $U_2$ (C5) & $-\frac i4$ printed; correct is $+\frac i4$ & none\\
\eqI{32}--\eqI{33} (S1) & overall sign of the 4-point contribution reversed & overall sign of the bare $M^2$\\
\eqI{39} (C6) & $\bar\eta_{\rho\sigma}\to\bar\eta^{\rho\sigma}$ & none\\
\eqI{40} (B1) & finite part $+2$ printed; correct is $-2$ & absorbable in counterterms\\
\eqI{47}--\eqI{49} (S1) & overall sign reversed in all nine terms & carried into the later calculation\\
\eqI{54}--\eqI{59} & sign reversed (inherited); in \eqI{59} $-\frac72\to-\frac52$; in the text ``4th''$\to$``2nd'' & sign: inherited (S1); \eqI{59}: absorbable in $\Delta\alpha_3$ (B2); ``4th'': typo (C7)\\
\eqI{60}--\eqI{62} & agreement & ---\\
\eqI{63}--\eqI{67} & sign reversed (inherited); no new error & ---\\
\textbf{Overall} (S1) & in the paper the 4-point and 3-point contributions are consistent with each other, both reversed; the local cancellation \eqI{44}--\eqI{67} holds in both conventions & bare $M^2=-$(definition)\\
\bottomrule
\end{tabular}
\end{center}
\end{minipage}


% ---------------------------------------------------------------------
\subsection{Code checks \texorpdfstring{\tC}{}}
\label{sec:ch04:t10code}
% ---------------------------------------------------------------------

The script \texttt{code/04\_three\_point\_local\_log/check\_60\_67.py} (SymPy; about 8\,s) checks the following.
It was actually run.
It is a check only and does not replace the analytic derivation above.
\begin{itemize}[leftmargin=1.2em,itemsep=3pt]
\item[(a)] \textbf{\eqI{62}.} The first three terms of \eqI{14}, expanded at $D=4+\epsilon$, give $\frac{H^2}{16\pi^2}\{4/y-2\ln(y/4)-1\}$.
The $n$-th series term vanishes identically at $D=4$.
The printed $c_0$ makes a difference only at order $D-4$.
\item[(b)] \eqI{62} satisfies $(4y-y^2)A''+4(2-y)A'=3k$, with $k=2h$.
\item[(c)] \textbf{Lemma (D5)}, verified on a test function with general $D$.
In the script, $\ln(aa')\times[\text{distribution}]$ acting on $g$ is represented as $l\,\mathcal O(g)+\mathcal O(l\,g)$, since $l'=\ln a(\eta')$ multiplies $g(x')$ inside the integral.
In addition, (D3) and $\partial_0\big(a^p\partial_0\ln a\big)=(D-1)H^2a^D$ were verified.
\item[(d)] \textbf{Pointwise check of the non-local parts.}
At $D=4$, $\mathcal V_1$, $\mathcal W$, $\mathcal V_3$ and $\mathcal V_4$ were differentiated by brute force.
This used explicit coordinates, $i\Delta_A=A+k\ln(aa')$, and the actual $A(y)$ of \eqI{62}.
At non-coincident points the delta-function terms vanish.
The results agree exactly with the non-local parts of \eqref{eq:ch04:63}--\eqref{eq:ch04:V4} at three random points, with 30-digit precision.
\item[(e)] \textbf{General-$D$ sum.} The coefficients are $\frac{D-1}{D-3}$, $-\frac4{D-3}$, $\frac{(D-1)^2}{D-3}$, $-\frac{2(D-4)}{D-3}$ and $\frac2{D-3}$.
Together with \eqI{44}, the sum is $-\frac{(D-4)(D+1)}{4(D-3)}[\ldots]$.
\item[(f)] The paper's \eqI{63}--\eqI{66} are $-1$ times the derived terms, term by term, and their sum gives the paper's \eqI{67}.
\end{itemize}

\paragraph{Actual output} (Python 3, SymPy 1.14.0; long lines wrapped, marked by \texttt{\textbackslash}).
The output of the re-run in the code folder is identical to the output in the original log.
In section (e) the symbols mean (in units of $\kappa^2k$):
\texttt{LB}~$=L\,a^D\square\,i\delta^D$, \texttt{LL}~$=L\,a^{D-2}\nabla^2i\delta^D$, \texttt{M}~$=H^2a^Di\delta^D$,
\texttt{Q4}~$=(aa')^{D-2}L\nabla^4A$, \texttt{T}~$=H(aa')^{D-2}a\partial_0\nabla^2A$, \texttt{Tp}~$=H(aa')^{D-2}a'\partial'_0\nabla^2A$.
In section (f) the same symbols are in units of $\kappa^2H^2/(8\pi^2)$ at $D=4$.
{\scriptsize
\begin{verbatim}
== (a) Eq. (62): A(y) at D = 4 + eps from (14) with c0 = Gamma(D/2-1) ==
  first three terms -> H^2/(16 pi^2){4/y - 2 ln(y/4) - 1} : True
  n-th series term at D=4 vanishes identically             : True
  printed c0 changes A only at O(D-4)                      : True
== (b) A(y) of (62) solves (4y-y^2)A'' + 4(2-y)A' = 3k with k = 2h  (Box A = (D-1) k H^2) ==
   True
== (c) delta identities on a test function g, general D (l = ln a, p = D-2) ==
  -2i d0[a^p l d0 delta] == i ln(aa')(a^D Box - a^{D-2}Lap) delta + i(D-1)H^2 a^D delta : True
  ln(aa') Lap delta acting on g == 2 l Lap g (time-only factor)        : True
  d0(a^p d0 ln a) = (D-1) H^2 a^D                                      : True
== (d) pointwise check of the non-local parts, D = 4, actual A(y) of (62), k = 2h ==
  point 0: |V1-R1|,|W-RW|,|V3-R3|,|V4-R4| = [0.0, 0.0, 0.0, 0.0]  \
           (scales [0.0553, 0.0515, 0.0639, 0.0428] )
  point 1: |V1-R1|,|W-RW|,|V3-R3|,|V4-R4| = [0.0, 0.0, 0.0, 0.0]  \
           (scales [0.00320, 0.00321, 0.00353, 0.00289] )
  point 2: |V1-R1|,|W-RW|,|V3-R3|,|V4-R4| = [0.0, 0.0, 0.0, 0.0]  \
           (scales [0.00203, 0.00597, 0.00544, 0.00255] )
== (e) general-D assembly, units kappa^2 k: -V1 + (D-1)/(D-3) W - V3 - V4 ==
  coefficients: LB: (D - 1)/(D - 3)  LL: -4/(D - 3)  M: (D - 1)**2/(D - 3)  Q4: -2*(D - 4)/(D - 3)  \
                T: 2/(D - 3)  Tp: 2/(D - 3)
  derived (44)_log-local + derived (67)_local = \
                -(D - 4)*(D + 1)*(D**2*M + D*LB - 2*D*M - LB - 4*LL + M)/(4*(D - 3))
  at D=4: LB, LL, M coefficients of (67) = [3, -4, 9]  T,Tp: [2, 2]
== (f) the paper's (63)-(66) and (67), units kappa^2 H^2/(8 pi^2) at D=4 ==
  each equals minus the derived term: True True True True
  paper sum: -3*LB + 4*LL - 9*M - 2*T - 2*Tp    (paper's (67): -3LB + 4LL - 9M - 2T - 2Tp)
\end{verbatim}
}
\noindent
The factor in (e) is $D^2M+D\,LB-2DM-LB-4LL+M=(D-1)LB-4LL+(D-1)^2M$, which is the bracket of \eqref{eq:ch04:cancel}.

% ---------------------------------------------------------------------
\subsection{Calculation record}
% ---------------------------------------------------------------------
\begin{itemize}[leftmargin=1.2em,itemsep=2pt]
\item \textbf{Derived.} \eqI{60}--\eqI{61} [Eqs.~\eqref{eq:ch04:60}--\eqref{eq:ch04:61}]; \eqI{62} [Eq.~\eqref{eq:ch04:62}];
distributional tools (D1)--(D5); \eqI{63}--\eqI{66} at general $D$ [Eqs.~\eqref{eq:ch04:63}--\eqref{eq:ch04:V4}];
\eqI{67} at general $D$ [Eq.~\eqref{eq:ch04:67D}] and at $D=4$ [Eq.~\eqref{eq:ch04:67}];
justification of the passage to $D=4$ (Section~\ref{sec:ch04:Dto4});
cancellation against \eqI{44} [Eq.~\eqref{eq:ch04:cancel}].
\item \textbf{Checks passed.} SymPy (a)--(f), actually run.
\item \textbf{Discrepancies.} \eqI{60}--\eqI{62} agree.
In \eqI{63}--\eqI{67} the overall sign is reversed. This is inherited from \eqI{47}; there is no new error.
\item \textbf{Settled.} The local terms of \eqI{67} and the $\ln(aa')$ and $H^2$ local terms of \eqI{44} cancel each other at general $D$, up to a remainder of order $D-4$.
The relative sign is consistent.
\item \textbf{Open point (at the time of Task 10).} The normal contributions (\eqI{68} and later), and the treatment of the term $2H(aa')^2(a\partial_0+a'\partial'_0)\nabla^2A$ of \eqI{67}.
\end{itemize}
